\section{The line graph in dimension six}
\label{sec:line}
\begin{lemma}\label{lem:weighted}
An independent set in $\Gamma_6$ containing $t$ even-weight vectors and $o$
odd-weight vectors satisfies $3t+4o\leq19$. Its cardinality is at most six;
cardinality six requires $(t,o)=(5,1)$.
\end{lemma}
\begin{proof}
Write $j=(1,1,1,1,1,1)$. The even-weight vectors lie in $j^\perp$, whose
restricted form has radical $\Span(j)$ and rank four. If $o=0$, the Gram
matrix of the selected vectors is $J_t+I_t$, of rank $t$ for even $t$ and
$t-1$ for odd $t$. Thus $t\leq5$.

If $o>0$, fix an odd selected vector $a$. For each even selected vector $e$,
the vector $z_e=e+a$ belongs to $a^\perp$. These vectors are orthonormal and
span a nondegenerate $t$-space $A$. The ambient space $a^\perp$ is
nondegenerate of dimension five. The span $D$ of the vectors $u+a$, over odd
selected vectors $u$, is totally isotropic and orthogonal to $A$ inside
$a^\perp$. Hence $2\dim D\leq5-t$. The odd selected vectors are distinct
members of $a+D$, so $o\leq2^{\dim D}$. Consequently $o\leq4$ for $t=0,1$,
$o\leq2$ for $t=2,3$, and $o\leq1$ for $t=4,5$. These bounds prove both claims.
\end{proof}

\begin{theorem}\label{thm:lines}
$12\leq\chi(\Gamma_6)\leq13$.
\end{theorem}
\begin{proof}
There are $31$ nonzero even-weight and $32$ odd-weight vectors. Their total
weight, using weights three and four respectively, is $221$. By
Lemma~\ref{lem:weighted}, every color class has weight at most $19$, giving
$\chi(\Gamma_6)\geq\lceil221/19\rceil=12$.
For the upper bound, interpret an integer $x$ as the vector whose coordinate
$i$ is its $i$th binary digit. The certificate \texttt{results/line-coloring.json}
partitions the $63$ nonzero vectors into $13$ classes. The independent checker
verifies that every pair of distinct vectors in each class has dot product one.
\end{proof}

For completeness, the thirteen color classes, in the same binary integer
coordinates, are displayed below. They partition $\{1,\ldots,63\}$, and
their within-class dot products can be checked directly.
\[
\begin{gathered}
\{3,29,45,53,57,62\},\quad \{12,20,39,56,59\},\quad
\{15,18,43,52,61\},\\
\{23,24,44,47,48\},\quad \{16,50,51,54,58\},\quad
\{32,46,55,60\},\\
\{13,22,42,49,63\},\quad \{5,19,28,38,41\},\quad
\{8,10,25,40\},\\
\{2,26,30,35\},\quad \{1,7,9,17,33\},\quad
\{6,11,27,34,37\},\quad \{4,14,21,31,36\}.
\end{gathered}
\]

Any hypothetical $12$-coloring has a class of size six. By the lemma it
consists of an odd vector $a$ and five even vectors $e$. The vectors $a$ and
$e+a$ form an orthonormal basis. An isometry therefore maps this class to
$\{1,3,5,9,17,33\}$. Fixing this class reduces the remaining question to an
$11$-coloring of the induced graph on the other $57$ vectors.
